\chapter{Introduzione e contesto}
\label{cap1}
\section{Introduzione e motivazioni}
Nelle proprie interazioni, l'essere umano ha la capacità di anticipare il comportamento degli altri individui e di agire di conseguenza. Questa abilità di anticipare un'azione è intrinseca della nostra specie ed è involontaria, e ne facciamo largo uso in quasi tutte le interazioni umane. Lo studio di questa capacità all'interno del campo della computer vision prende il nome di Human Activity Analysis (HAA), la quale presenta una vasta gamma di applicazioni, tra le quali figurano la rilevazione di comportamenti anomali nell'ambito della videosorveglianza, l'avviso di pericolo nei sistemi avanzati di assistenza alla guida, la realtà aumentata e il recupero video.
Nell'ambito della robotica, un robot collaborativo che lavora al fianco di un operatore umano necessita di anticipare le intenzioni di quest'ultimo per coordinare i propri movimento, evitare collisioni e offrire supporto nel momento esatto. In aggiunta, nell'ambito della realtà aumentata, un sistema indossabile che vuole sovrapporre informazioni contestuali utili deve prevedere quale oggetto l'utente sta per usare prima che l'azione inizi.  
È chiaro come il comune denominatore di questi svariati impieghi risieda nella comprensione video, ovvero ricavare informazioni dall'osservazione delle sequenze di azioni. Il valore di tali informazioni, però, decade rapidamente nel tempo. Sapere che un pedone attraverserà la strada è utile se questa informazione viene ottenuta con almeno qualche secondo di anticipo; è inutile invece se arriva dopo che l'evento si è già verificato. Questa caratteristica rende l'anticipazione di un'azione sostanzialmente diversa dal suo riconoscimento. 

\section{Stato dell'arte}
Per comprendere la specificità dell'action anticipation occorre collocarla rispetto ai task correlati nell'ampio panorama della comprensione video, facendo distinzione tra action recognition, action prediction (o early action recognition) e action anticipation. 
Per action recognition si intende la classificazione di un'azione umana basandosi sulla sua completa esecuzione.
L'action prediction, che spesso viene chiamata anche early recognition, mira a prevedere un'azione che è attualmente in corso, a partire solamente da un'osservazione parziale. Il suo obiettivo è riconoscere l'azione il prima possibile, mentre è ancora in fase di svolgimento. 
Per action anticipation, invece, si intende la previsione di un'azione ancor prima che essa abbia avuto inizio. L'osservazione termina prima che l'azione futura abbia inizio. Questa distinzione non è banale in quanto, nell'action recognition, il modello ha accesso tutti i componenti che caratterizzano l'azione: la fase iniziale, lo svolgimento centrale e la fase conclusiva, mentre nell'early recognition ha accesso almeno alla fase iniziale, contrariamente allo scenario dell'anticipation, dove il modello deve prevedere un evento completamente nuovo sfruttando una comprensione profonda di segnali indiretti, ovvero le azioni precedenti, gli oggetti presenti nella scena, la postura del corpo o la direzione dello sguardo, tutti elementi che costituiscono il contesto da cui dedurre l'intenzione futura. Questa differenza ha profonde implicazioni sulle differenze di difficoltà dei task: l'anticipation è intrinsecamente più incerta della recognition poiché, dall'osservazione di uno stesso evento si ricavano più azioni future plausibili.

\section{Obiettivo e contributo del lavoro}
